\chapter{Conclusões e Recomendações}
\label{cap:05}

\section{Conclusões}

O objetivo geral deste trabalho foi desenvolver o GUY, uma extensão para o \textit{Visual Studio Code} destinada à análise estrutural de código-fonte e à geração automática do Grafo de Fluxo de Controle (GFC). Dentro do escopo definido, a ferramenta implementada contemplou a geração visual do GFC, o cálculo da Complexidade Ciclomática \(V(G)\), a navegação entre grafo e código-fonte e a execução local da análise. Os resultados obtidos indicam que a solução atende à proposta de aproximar a estrutura lógica do código de uma representação visual capaz de apoiar sua compreensão e o planejamento de testes estruturais.

A análise de ferramentas como SonarQube, Codacy e ESLint permitiu delimitar a proposta do GUY. Enquanto essas soluções concentram seus resultados em painéis, diagnósticos textuais e regras de qualidade, o GUY foi concebido com foco na visualização interativa do fluxo de controle e na relação entre os elementos do grafo, a métrica de complexidade e o código-fonte analisado.

A definição dos requisitos funcionais e não funcionais orientou a prototipagem e a implementação da extensão. Foram considerados, entre outros aspectos, a execução local, a integração com o VS Code, a geração do GFC, o cálculo de \(V(G)\), a indicação de caminhos independentes e a interação entre grafo e código-fonte. Esses requisitos também serviram como referência para a avaliação do protótipo desenvolvido.

A implementação foi realizada em \textit{TypeScript}, com análise sintática de Python por Tree-sitter em WebAssembly, construção do GFC, cálculo de métricas estruturais e renderização interativa em uma Webview desenvolvida com React e \texttt{@xyflow/react}. O resultado foi um protótipo funcional organizado em módulos responsáveis pela análise do código, construção do grafo, cálculo de métricas, navegação no editor e interface visual.

A avaliação foi realizada com três arquivos de exemplo em Python, escolhidos para representar estruturas fundamentais de controle de fluxo. Nos cenários analisados, os grafos produzidos apresentaram relações coerentes entre nós, arestas e decisões, e os valores de Complexidade Ciclomática foram compatíveis com a expressão \(V(G)=E-N+2P\). O primeiro experimento resultou em \(V(G)=4\), o segundo em \(V(G)=2\) e o terceiro em \(V(G)=3\), valores também correspondentes à quantidade de decisões acrescida de uma unidade nos grafos conexos avaliados.

Esses resultados demonstram consistência nos cenários selecionados, mas não constituem uma validação abrangente da ferramenta. A amostra foi restrita a três exemplos controlados, sem incluir projetos reais de maior porte, comparação sistemática com uma implementação de referência ou avaliação com usuários. Dessa forma, as conclusões devem permanecer limitadas ao escopo e aos casos analisados neste trabalho.

Do ponto de vista qualitativo, destacam-se a integração direta ao VS Code, a execução local, a visualização interativa do GFC, o relacionamento entre elementos do grafo e trechos do código-fonte, os modos de visualização simplificado e detalhado e a possibilidade de analisar um arquivo completo, uma seleção ou a função atual. Essas características indicam o potencial da ferramenta para apoiar atividades educacionais e análises estruturais dentro do escopo avaliado, especialmente por aproximarem a métrica \(V(G)\) da estrutura visual do código.

As limitações observadas delimitam o alcance da solução. A versão atual concentra-se em Python e em estruturas de controle comuns, sem representar de forma completa exceções, chamadas entre funções, recursão, geradores, programação assíncrona e efeitos dinâmicos de execução. Além disso, por se tratar de análise estática, a ferramenta não substitui a execução de testes nem garante a identificação de todos os comportamentos possíveis em tempo de execução. Há também limitações de visualização em grafos extensos, que podem se tornar densos em programas com muitas estruturas aninhadas, mesmo com o uso do modo simplificado e de mecanismos para reduzir a quantidade de informações exibidas.

\section{Trabalhos Futuros}

Como evolução do analisador, recomenda-se ampliar o suporte a recursos da linguagem Python, incluindo exceções, chamadas entre funções, recursão, geradores e programação assíncrona, preservando a coerência entre os elementos sintáticos reconhecidos e a construção do GFC. Também pode ser considerada a ampliação da ferramenta para outras linguagens de programação por meio de novas gramáticas do Tree-sitter.

Para fortalecer a avaliação técnica, recomenda-se criar uma suíte de testes automatizados baseada em exemplos de referência, com grafos e métricas esperados definidos previamente. Essa abordagem permitiria comparar continuamente os resultados produzidos pela ferramenta durante sua evolução e reduzir o risco de regressões na construção do GFC e no cálculo da Complexidade Ciclomática.

Também se recomenda ampliar a validação com projetos reais de diferentes tamanhos e níveis de complexidade, além de realizar estudos com estudantes e desenvolvedores. Uma avaliação com usuários permitiria investigar aspectos como usabilidade, compreensão dos grafos, clareza da interação entre grafo e código-fonte e utilidade da ferramenta no planejamento de testes.

Por fim, recomenda-se aprimorar os recursos de navegação e organização visual para grafos extensos, incluindo estratégias de agrupamento, filtragem, expansão progressiva e exportação da visualização. Esses recursos podem reduzir a densidade visual e tornar a análise mais adequada a programas de maior porte.
